% Original PDF: Section 2, pages 2--5.
\section{Preliminaries and the MCES algorithm}\label{sec:preliminaries}
\subsection{Finite discounted and proper episodic MDPs}
Let $\mathcal S$ be a finite set of nonterminal states, with finite nonempty action sets $\mathcal A(s)$, and put
$\mathcal I=\{(s,a):s\in\mathcal S,\ a\in\mathcal A(s)\}$ and $J=|\mathcal I|$. A transition from $(s,a)$ produces a reward and
a next state according to the MDP kernel. Expected absolute one-step rewards are finite.
Terminal states, when present, have value zero and are omitted from $\mathcal I$. Write $\mathcal P$ for the
finite set of deterministic stationary policies.

The discount factor is $\gamma\in[0,1]$. If $\gamma=1$, assume that every $\pi\in\mathcal P$ is \emph{proper}: from
every nonterminal state it reaches a terminal state almost surely. If $\gamma<1$, no termination
assumption is imposed. State and action values are
\[
v_\pi(s)=\mathbb E_\pi\!\left[\sum_{t\ge0}\gamma^tR_t\mid S_0=s\right],\qquad
q_\pi(s,a)=\mathbb E_\pi\!\left[\sum_{t\ge0}\gamma^tR_t\mid S_0=s,A_0=a\right],
\]
where rewards after termination are zero and the initial action in the second expression is
prescribed. The policy is followed thereafter. The values are finite: discounting suffices when
$\gamma<1$, while properness makes each finite nonterminal transition matrix transient when
$\gamma=1$. In particular,
\begin{equation}\label{eq:bound}
B:=\max_{\pi\in\mathcal P}\|q_\pi\|_\infty<\infty.
\end{equation}
All vector inequalities below are coordinatewise, and all norms are supremum norms unless
indicated otherwise. We use $\inf\varnothing=\infty$.

\begin{lemma}[Uniform termination in the episodic case]\label{lem:termination}
If every deterministic stationary policy is proper, there are an integer $m\ge1$ and $\rho\in[0,1)$ such that the termination
time $T$ satisfies $\mathbb P(T>jm)\le\rho^j$ under every control law, uniformly over starting states and
prescribed initial actions. In particular, $\sup\mathbb E[T^2]<\infty$.
\end{lemma}
\begin{proof}
Let $w_n(s)$ be the largest probability of remaining nonterminal after $n$ transitions,
allowing arbitrary history-dependent policies. Then $w_0(s)=1$ and finite-horizon dynamic
programming gives
\[
w_{n+1}(s)=\max_{a\in\mathcal A(s)}\sum_{s'\in\mathcal S}P(s'\mid s,a)w_n(s').
\]
The sequence decreases to a vector $w$. Finiteness permits passage through the maximum
and sum. A deterministic stationary maximizer $\pi$ therefore satisfies $w=P_\pi w$. Properness
implies $P_\pi^n\to0$, so $w=0$. Choose $m$ with $\max_s w_m(s)\le\rho<1$. Conditioning at successive
blocks of $m$ transitions gives the asserted bound, also when the first action is prescribed.
Summing the geometric tail gives a uniform second moment.
\end{proof}

For a deterministic policy $\rho$, write $T_\rho v=r_\rho+\gamma P_\rho v$, and let $Tv(s)=\max_a\{r(s,a)+
\gamma\sum_{s'}P(s'\mid s,a)v(s')\}$. The following form of policy improvement will be used repeatedly;
see also Howard (1960).

\begin{lemma}[Policy improvement and Bellman verification]\label{lem:improvement}
Suppose that, for every
state,
\[
q_\pi(s,\rho(s))\ge q_\pi(s,\pi(s)).
\]
Then $v_\rho\ge v_\pi$ and $q_\rho\ge q_\pi$. If additionally $q_\rho\ne q_\pi$, then $\sum_{i\in\mathcal I}q_\rho(i)>\sum_{i\in\mathcal I}q_\pi(i)$. A
deterministic policy greedy for its own action values is optimal, including among history-dependent policies.
\end{lemma}
\begin{proof}
The premise is $T_\rho v_\pi\ge v_\pi$. Since $v_\rho=T_\rho v_\rho$,
\begin{equation}\label{eq:inverse}
v_\rho-v_\pi=(I-\gamma P_\rho)^{-1}(T_\rho v_\pi-v_\pi)\ge0.
\end{equation}
The inverse is the nonnegative series $\sum_{j\ge0}(\gamma P_\rho)^j$: discounting or properness ensures convergence. The identity $q_\pi(s,a)=r(s,a)+\gamma\sum_{s'}P(s'\mid s,a)v_\pi(s')$ gives the action-value
comparison. Distinct coordinatewise ordered vectors have strictly ordered sums.

If $\pi$ is greedy for $q_\pi$, then $Tv_\pi=v_\pi$. Iterating the Bellman inequality along any control
law shows that its expected return is at most $v_\pi$, up to a terminal remainder bounded by
$\|v_\pi\|_\infty\gamma^n$ when $\gamma<1$, or by $\|v_\pi\|_\infty\mathbb P(T>n)$ when $\gamma=1$. That remainder tends to zero, in
the latter case by Lemma~\ref{lem:termination}. The policy $\pi$ attains $v_\pi$, proving optimality.
\end{proof}

An optimal deterministic stationary policy exists. For example, choose $\pi$ maximizing
$\sum_s v_\pi(s)$ over the finite set $\mathcal P$. If it were not greedy for $q_\pi$, a greedy improvement $\rho$ would
have $T_\rho v_\pi-v_\pi\ge0$ with a positive coordinate; the identity term in the inverse in \eqref{eq:inverse} would
give a strict increase of the sum, a contradiction. We denote the optimal state and action
values by $v^\star$ and $q^\star$, respectively. We do not use the converse implication from action-value
ordering to state-value ordering. Equal action-value targets can belong to policies with
different state values.

\subsection{Initial-visit Monte Carlo control}\label{sec:algorithm}
Let $q_k\in\mathbb R^{\mathcal I}$ be the current estimate, starting from a deterministic finite $q_0$. At the beginning
of iteration $k$, select a deterministic greedy policy
\begin{equation}\label{eq:greedy}
\pi_k\in\mathcal G(q_k),\qquad
\mathcal G(q)=\{\pi\in\mathcal P:q(s,\pi(s))=\max_a q(s,a)\text{ for all }s\}.
\end{equation}
Let $\mathcal H_k$ contain the history after computing $q_k$ but before this selection. We allow either of
two tie rules, fixed throughout a run:
\begin{enumerate}
\item \emph{Fixed priorities.} At every state use a fixed priority ordering of the actions, including at
initialization.
\item \emph{Full inertia.} Choose $\pi_0$ greedily, and thereafter retain the preceding action whenever it
remains greedy:
\begin{equation}\label{eq:inertia}
\pi_{k+1}(s)=\pi_k(s)\quad\text{if }\pi_k(s)\in\arg\max_a q_{k+1}(s,a).
\end{equation}
Otherwise choose any maximizing action, possibly at random and using the history.
No statewise independence is required for this rule.
\end{enumerate}
After selecting $\pi_k$, choose the initial-pair distribution $\sigma_k$ and let $\mathcal F_k$ contain $\mathcal H_k$, $\pi_k$, and $\sigma_k$.
The distribution may depend on the policy and on the entire past. Conditional on $\mathcal F_k$, draw
$I_{k+1}$ with distribution $\sigma_k$ and obtain a return estimate $G_{k+1}$. In the episodic implementation,
start with that pair, follow $\pi_k$ thereafter until termination, and sum the discounted rewards.
Only the initial pair is updated:
\begin{equation}\label{eq:update}
q_{k+1}(i)=q_k(i)+\alpha_k\mathbf1_{\{I_{k+1}=i\}}(G_{k+1}-q_k(i)),\qquad i\in\mathcal I.
\end{equation}
The policy remains fixed during each episode and is selected again after this update. The
next pre-policy history contains the pair, return, and update; in particular,
\[
\mathcal H_k\subseteq\mathcal F_k\subseteq\sigma(\mathcal F_k,I_{k+1})\subseteq\mathcal H_{k+1}\subseteq\mathcal F_{k+1}.
\]
All stopping times in the convergence proof refer to $(\mathcal F_k)$. This is the initial-visit version
of MC-O-PI, also called Monte Carlo control with exploring starts (MCES). First-visit and
every-visit algorithms update additional pairs within an episode; those algorithms are not
covered by the convergence theorem below.

\begin{assumption}[Cumulative learning, returns, and step sizes]\label{asp:main}
The following conditions hold:
\begin{enumerate}[label=(\arabic*)]
\item Almost surely, initial actions are uniform within each state, and every state accumulates
infinite learning:
\begin{equation}\label{eq:learning}
\begin{gathered}
\sigma_k(s,a)=\sigma_k(s)\ge0,\qquad
\sum_s|\mathcal A(s)|\sigma_k(s)=1,\qquad k\ge0,\\
\sum_{k=0}^{\infty}\alpha_k\sigma_k(s)=\infty\quad\text{for every }s.
\end{gathered}
\end{equation}
Here $\sigma_k(s)$ is the probability of each pair at state $s$, not its state marginal, which is
$|\mathcal A(s)|\sigma_k(s)$.
\item For a deterministic $c_v<\infty$, the return is conditionally unbiased with uniformly bounded
conditional variance. On pairs with positive conditional sampling probability,
\begin{equation}\label{eq:moments}
\begin{aligned}
\mathbb E[G_{k+1}\mid\mathcal F_k,I_{k+1}=i]&=q_{\pi_k}(i),\\
\mathbb E[(G_{k+1}-q_{\pi_k}(i))^2\mid\mathcal F_k,I_{k+1}=i]&\le c_v.
\end{aligned}
\end{equation}
No return law is required on a zero-probability initial-pair event.
\item The deterministic steps satisfy
\begin{equation}\label{eq:steps}
0<\alpha_k\le1,\qquad\sum_{k=0}^{\infty}\alpha_k^2<\infty,\qquad
\alpha_{k+1}\le\alpha_k\quad(k\ge k_{\mathrm m}),
\end{equation}
for a deterministic integer $k_{\mathrm m}$. Thus nonincrease is required only eventually. Equation~\eqref{eq:learning} already implies $\sum_k\alpha_k=\infty$.
\end{enumerate}
\end{assumption}

The state law can be unequal, adaptive, and zero for arbitrarily long intervals. Equivalently, draw a state with probability $d_k(s)$ and then choose its initial action uniformly, giving
$\sigma_k(s)=d_k(s)/|\mathcal A(s)|$. No lower bound or bound on waiting times is imposed in Assumption~\ref{asp:main}.
Infinite visitation alone is not a substitute for infinite \emph{learning weight} when a common global
step sequence is used.

Bounded one-step rewards suffice for \eqref{eq:moments} with ordinary episode returns. For $|R_t|\le R_{\max}$
and $\gamma<1$, the return is bounded by $R_{\max}/(1-\gamma)$. For $\gamma=1$ it is bounded by $R_{\max}T$, which
has a uniform second moment by Lemma~\ref{lem:termination}. Unbounded rewards or other estimators are
allowed when the explicit conditional moments hold. In a continuing discounted problem,
an infinite rollout is unnecessary: with bounded rewards, draw an independent $U$ with
$\mathbb P(U=j)=(1-\gamma)\gamma^j$, simulate from the prescribed initial pair using $\pi_k$ thereafter, and
return $R_U/(1-\gamma)$. This bounded estimator is unbiased; for $\gamma=0$ take $U=0$.

No comparability condition such as $\alpha_{k+j}\ge c\alpha_k$ for $j$ up to order $1/\alpha_k$ is assumed.
Nonincrease compares earlier seed updates with a state's eventual completion step and
controls future variance on that state's learning clock.
